\section{An integral cyclotomic class}
\label{sec:arithmetic}

The coefficient exact sequence identifies the obstruction to lifting a
degree-five coefficient class: it is a class in integral $K_4$ killed by
five.  We will show that this obstruction vanishes for one cyclotomic
ring.  The proof first uses localization to detect coefficient $K_4$
geometrically, then uses finite generation to eliminate all five-primary
torsion.  We finish by recording how the chosen integral lift behaves
under the four cyclotomic automorphisms.

Fix a primitive fifth root of unity $\zeta$, and put
\[
 N=\Q(\zeta),\qquad S=\Z[\zeta,1/5],\qquad
 u=\zeta-1,\qquad \Delta=\Gal(N/\Q).
\]
Let $\chi:\Delta\to\F_5^\times$ be the character defined by
$\delta(\zeta)=\zeta^{\chi(\delta)}$.  We use the coefficient convention and product of Section~\ref{sec:coefficients}.

\paragraph{The cyclotomic ring.}
We record the arithmetic facts that will control the localization
sequence.  Let $\Phi_5(X)=X^4+X^3+X^2+X+1$ be the fifth cyclotomic
polynomial.  The polynomial of $u$ is
\[
 \Phi_5(1+T)=T^4+5T^3+10T^2+10T+5,
\]
which is Eisenstein at five.  In particular it is irreducible over
$\Q_5$, so $N$ has exactly one place above five.  The ring $\Z_5[u]$ is
finite free over $\Z_5$; irreducibility makes it a domain, and integrality
gives dimension one.  Its reduction modulo five is $\F_5[T]/(T^4)$,
with unique maximal ideal generated by $T$.  Thus $\Z_5[u]$ is a
Noetherian local domain with maximal ideal $(5,u)$.  The defining equation
gives
\[
 u^4=-5(1+2u+2u^2+u^3).
\]
The factor in parentheses is a unit, since its residue modulo $(5,u)$
is one.  Hence $(5,u)=(u)$, and $\Z_5[u]$ is a discrete valuation ring.
Being integrally closed with fraction field $\Q_5(u)$, it is the integer
ring of that field, with residue field $\F_5$ and ramification index four.

The norm of $u$ is $\Phi_5(1)=5$.  Since
$\Phi_5'(\zeta)=5\zeta^4/(\zeta-1)$, the discriminant of the power
basis $1,\zeta,\zeta^2,\zeta^3$ is
\[
 \operatorname{N}_{N/\Q}\bigl(\Phi_5'(\zeta)\bigr)
 =\frac{5^4}{\operatorname{N}_{N/\Q}(u)}=125.
\]
Let $\mathcal O_N$ be the ring of integers of $N$.  The index of
$\Z[\zeta]$ in $\mathcal O_N$ can therefore have
no prime divisor other than five, by the discriminant-index formula.
The local calculation above shows that this order is already maximal
at five.  Thus the index is one:
\[
 \mathcal O_N=\Z[\zeta],\qquad \operatorname{disc}(N)=125.
\]
There are two complex places and no real places.  Minkowski's ideal-class
bound \cite[Chapter~4, Theorem~4.3]{MilneANT2020} is consequently
\[
 \frac{4!}{4^4}\left(\frac4\pi\right)^2\sqrt{125}<2.
\]
Every ideal class has an integral ideal representative of positive
integral norm at most this bound, hence of norm one.  It follows that
$\Pic(\mathcal O_N)=0$, and localization gives $\Pic(S)=0$.

The norm calculation also shows that $u$ and all its conjugates are
units in $S$.  Every closed point $v$ of $\Spec S$ has a finite
residue field $k_v$ of characteristic different from five.  The image of
$\zeta$ in $k_v$ has order exactly five: its order divides five, and it
cannot be one because $u$ is a unit.  These are precisely the residue
fields to which we will apply the finite-field clause of
Lemma~\ref{lem:bott}.

\paragraph{An equivariant coefficient class.}
The determinant splits the unit homomorphism
$S^\times\to K_1(S)$, so $[\zeta]\in K_1(S)$ has order
five.  The coefficient exact sequence supplies
$\beta'\in K_2(S;5)$ with
$\partial_{\mathrm{coef}}\beta'=[\zeta]$.  This coefficient group is an
$\F_5$-vector space by Lemma~\ref{lem:moore-product}, and $\Delta$ acts
on it by functoriality.  Define
\[
 e_\chi=\frac14\sum_{\delta\in\Delta}\chi(\delta)^{-1}\delta
       \in\F_5[\Delta],\qquad \beta=e_\chi\beta'.
\]
Here $1/4$ is the inverse of $4$ in $\F_5$.  Naturality makes the
coefficient boundary $\Delta$-equivariant, and
$\delta[\zeta]=\chi(\delta)[\zeta]$.  Hence
$\partial_{\mathrm{coef}}\beta=e_\chi[\zeta]=[\zeta]$.
Relabeling the sum gives $\tau e_\chi=\chi(\tau)e_\chi$ for every
$\tau\in\Delta$, so
\begin{equation}
 \label{eq:equivariant-bott}
 \partial_{\mathrm{coef}}\beta=[\zeta],\qquad
 \delta\beta=\chi(\delta)\beta\quad(\delta\in\Delta).
\end{equation}
For any $S$-scheme on which the operation is used, let $b$ denote
left multiplication by the restriction of $\beta$, and put
$b^2=b\circ b$.  Thus $b^2(x)=\beta\cdot(\beta\cdot x)$, with the
indicated parentheses.  The coefficient localization boundaries below
are different from $\partial_{\mathrm{coef}}$; Lemma~\ref{lem:bott-localization}
states that $b$ commutes with the former.

The next lemma is the arithmetic reason that localization will leave
only a single coefficient line.  It concerns reductions of integral
$K_2$-classes, not all of $K_2(N;5)$.

\begin{lemma}[Unramified degree-two classes]
\label{lem:cyclotomic-unramified}
Let $N=\Q(\zeta)$ and $S=\Z[\zeta,1/5]$ as above.  If a class in
$K_2(N)/5$ has zero tame residue in $k_v^\times/k_v^{\times5}$ at every
closed point $v$ of $\Spec S$, then the class is zero.
\end{lemma}

\begin{proof}
Matsumoto's symbol theorem
\cite[\S5, Theorem~5.10 and Corollary~5.11]{Matsumoto1969} gives the
classical Steinberg-symbol presentation; see \cite[III.6.1]{Weibel2013}
for its modern formulation.  The classical-to-Quillen comparison
\cite[IV.1.10.1 and IV.7.2]{Weibel2013} identifies this group with
$K_2(N)$.  The degree-two norm-residue theorem
\cite[Theorem~6.16]{Voevodsky2011} then gives
\[
 K_2(N)/5\simeq H^2(N,\mu_5^{\otimes2}).
\]
Use $\zeta$ to identify one copy of $\mu_5$ with the constant module
$\Z/5$.  The Kummer sequence identifies the resulting group
$H^2(N,\mu_5)$ with $\Br(N)[5]$.  Write $a$ for the Brauer class
corresponding to the class in the lemma.

For a closed point $v$ of $\Spec S$, the residue characteristic is
not five.  The Quillen boundary on $K_2(N)$ is the tame symbol.  Under
norm residue and the chosen twist identification, its reduction modulo
five and the Brauer residue agree up to an overall sign depending on
boundary conventions; in particular, their vanishing is equivalent.
Indeed, write each field element as a power of a uniformizer times a
unit.  Both residues kill unit--unit symbols, and on a symbol consisting
of a unit $\epsilon$ and a uniformizer they give the class of
$\bar\epsilon$, up to the corresponding sign.  Bilinearity and the
Steinberg relations give the general case
\cite[V.6.1.2 and V.(6.6)]{Weibel2013}
\cite[Proposition~7.10]{Harari2023Brauer}.  Thus $a$ has zero Brauer
residue at every finite place $v\nmid5$.

The completion $N_v$ has finite residue field $k_v$.  For its henselian
valuation ring, the kernel of the Brauer residue is $\Br(k_v)=0$
\cite[Proposition~6.7 and Example~6.8(a)]{Harari2023Brauer}.
Consequently $a|_{N_v}=0$ at every $v\nmid5$.  The infinite completions
are complex and have zero Brauer group.  Only the single place above
five remains.  In the global Brauer invariant sequence, the sum of all
local invariants is zero, so the invariant at this last place is zero.
The local invariant is injective there, and the global restriction map
is injective.  These are the local and global Brauer statements in
\cite[IV, Example~2.14(c)--(e), Proposition~4.3 and Remark~4.4(a)]{MilneCFT2020}; see also
\cite[III.6 and VI.8.1.1]{Weibel2013}.  They imply $a=0$, hence the
original class is zero.  No tame-residue calculation at the place above
five is used.
\end{proof}

We can now eliminate the obstruction in integral $K_4$.  The first part
of the proof detects coefficient classes after passage to an algebraic
closure; the last part explains why this also rules out integral
five-primary torsion.

\begin{lemma}[The integral torsion obstruction vanishes]
\label{lem:cyclotomic-k4}
For $S=\Z[\zeta,1/5]$, one has $K_4(S)[5]=0$.
\end{lemma}

\begin{proof}
Because $S$ is a Dedekind ring with $\Pic(S)=0$, every divisor
on $\Spec S$ is principal.  Equivalently, the normalized valuation
map
\[
 N^\times\longrightarrow\bigoplus_{v\in(\Spec S)^{(1)}}\Z,
 \qquad f\longmapsto\bigl(v(f)\bigr)_v,
\]
is surjective.  The sums here and below have finite support.  Quillen
localization and d\'evissage for $S$
\cite[\S7, Proposition~3.2 and paragraph~3.4]{Quillen1973Higher} identify the
boundary from $K_1(N;5)$ with this valuation map modulo five.  We use
the positive normalized valuation convention
\cite[Example~V.6.1.2 and (V.6.6)]{Weibel2013}.  Thus
\[
 K_1(N;5)\longrightarrow\bigoplus_vK_0(k_v;5)
\]
is surjective.

The number field $N$ has $5$-cohomological dimension two
\cite[VI.\S7 and VI.8.2]{Weibel2013}.  The odd-degree clause of
Lemma~\ref{lem:bott} therefore makes
$b^2:K_1(N;5)\to K_5(N;5)$ an isomorphism.  For every $v$, the finite
field $k_v$ contains the primitive image of $\zeta$, so the finite-field
clause makes $b^2:K_0(k_v;5)\to K_4(k_v;5)$ an isomorphism as well.
The same restricted class $\beta$ acts on every term.  Its compatibility
with localization in Lemma~\ref{lem:bott-localization} carries the last
surjection to a surjection
\[
 K_5(N;5)\longrightarrow\bigoplus_vK_4(k_v;5).
\]
In the exact localization segment
\[
 K_5(N;5)\longrightarrow\bigoplus_vK_4(k_v;5)
 \longrightarrow K_4(S;5)\longrightarrow K_4(N;5),
\]
the middle arrow is consequently zero.  We obtain an injection
\begin{equation}
 \label{eq:k4-restriction}
 K_4(S;5)\hookrightarrow K_4(N;5).
\end{equation}

We next determine where this injected group can lie.  Take
$x\in K_4(S;5)$ and write $x_N$ for its image.  The even-degree
clause of Lemma~\ref{lem:bott} gives
$b:K_2(N;5)\xrightarrow{\sim}K_4(N;5)$, so $x_N=b(z)$ for some
$z\in K_2(N;5)$.  The localization residues of $x_N$ in
$K_3(k_v;5)$ vanish because $x$ extends over $S$.  Write
$\partial_v$ for the localization boundary at $v$.  Compatibility with
$b$ gives
\[
 0=\partial_v(x_N)=b\bigl(\partial_v z\bigr).
\]
Here $\partial_v z\in K_1(k_v;5)$, and
$b:K_1(k_v;5)\to K_3(k_v;5)$ is injective by the finite-field clause.
Thus every localization residue of $z$ is zero.

On the other hand,
$\partial_{\mathrm{coef}}z\in K_1(N)[5]=N^\times[5]=\langle\zeta\rangle$.
Choose $m\in\F_5$ such that this boundary equals $m[\zeta]$, and consider
the corrected class $z-m\beta_N$, where $\beta_N$ is the restriction of
$\beta$.  Equation~\eqref{eq:equivariant-bott} makes its coefficient
boundary zero.  Its localization residues are still zero, since $\beta$ extends
over $S$.  Exactness of the coefficient sequence identifies this
corrected class with an element of $K_2(N)/5$.  Naturality of reduction
with localization identifies its coefficient residues with the tame
residues of that element.  Lemma~\ref{lem:cyclotomic-unramified} therefore makes
the corrected class zero.  We have proved
\[
 z=m\beta_N,\qquad x_N=m\,b(\beta_N).
\]

Let $\overline N$ be an algebraic closure of $N$.  By the closed-field
clause of Lemma~\ref{lem:bott}, the restriction of $b(\beta_N)$ is
the nonzero square of a geometric Bott generator in
$K_4(\overline N;5)$.  Restriction is therefore injective on the
$\F_5$-line spanned by $b(\beta_N)$.  Together with
\eqref{eq:k4-restriction}, this proves
\begin{equation}
 \label{eq:k4-geometric-injection}
 K_4(S;5)\hookrightarrow K_4(\overline N;5).
\end{equation}

It remains to pass from coefficient detection to integral torsion.
The same closed-field clause gives $K_5(\overline N;5)=0$.  The
coefficient exact sequence then gives $K_4(\overline N)[5]=0$.
In fact $K_4(\overline N)$ has no nonzero element of any finite
five-power order: a suitable multiple of such an element would have
order five.

Set $G=K_4(S)$, and write $\rho$ for coefficient reduction at each
ring.  Naturality gives the commutative square
\[
 \begin{array}{ccc}
 G &\xrightarrow{\rho}& K_4(S;5)\\
 \big\downarrow && \big\downarrow\\
 K_4(\overline N)&\xrightarrow{\rho}&K_4(\overline N;5).
 \end{array}
\]
The right vertical arrow is injective by
\eqref{eq:k4-geometric-injection}.  If $g\in G$ has finite five-power
order, its integral image in $K_4(\overline N)$ is zero.  The square
therefore gives $\rho(g)=0$.  Since $G/5$ injects into
$K_4(S;5)$, this says $g\in5G$.  This conclusion holds for every
element of five-power order, not just for elements killed by five.

Finite generation now rules out a nonzero divisible primary subgroup.
Quillen's theorem makes $K_4(\mathcal O_N)$ finitely generated
\cite[\S1, Theorem~1 and Remark~(2)]{Quillen1973Finite}; see also
\cite[IV.6.9]{Weibel2013}.  Integral localization at the single place
above five has the segment
\[
 0=K_4(\F_5)\longrightarrow K_4(\mathcal O_N)
 \longrightarrow G\longrightarrow K_3(\F_5)=\Z/24.
\]
The two finite-field values are Quillen's integral computation
\cite{Quillen1972FiniteFields}; see \cite[IV.1.13]{Weibel2013}.
Thus $G$ is finitely generated.  Write
$G\simeq\Z^r\oplus P\oplus T$, where $P$ is its finite five-primary
torsion subgroup and $T$ is finite of order prime to five.  We proved
$P\subseteq5G$.  If $p=5g\in P$, the free component of $g$ is zero,
and its $T$-component is zero because multiplication by five is an
automorphism of $T$.  Hence $g\in P$, and $P\subseteq5P$.  The reverse
inclusion is automatic.  Iterating $P=5P$ to the exponent of the finite
group $P$ gives $P=0$.  In particular $K_4(S)[5]=0$.
\end{proof}

The coefficient obstruction has now been removed.  We choose the
degree-five class whose two Bott factors produce the square of the
cyclotomic character.

The cyclotomic-unit/Bott pattern has a methodological predecessor in
Soul\'e's coefficient and tower constructions
\cite[\S\S4.2--4.4, pp.~239--240]{Soule1987}.  The fixed integral
$K_5(S)$-lift below is a separate deduction from
$K_4(S)[5]=0$, proved here.

\begin{proposition}[Integral lift and conjugates]
\label{prop:integral-lift}
Let $u=\zeta-1\in S^\times$, and let $b$ be multiplication by the
class $\beta$ in \eqref{eq:equivariant-bott}.  There exists
$c\in K_5(S)$ such that
\begin{equation}
 \label{eq:integral-c}
 c\bmod5=b^2[u]_5\quad\text{in }K_5(S;5).
\end{equation}
For every $\delta\in\Delta$,
\begin{equation}
 \label{eq:conjugate-c}
 (\delta c)\bmod5=\chi(\delta)^2b^2[\delta u]_5
       \quad\text{in }K_5(S;5).
\end{equation}
For any fixed homomorphism $\iota:S\to H$ to a field, let
$\operatorname{res}_\iota$ denote extension of scalars.  Restriction of
these equalities gives the corresponding identities in $K_5(H;5)$.
In the second identity this means
\[
 \operatorname{res}_\iota(\delta c)\bmod5
 =\chi(\delta)^2 b_\iota^2[\iota(\delta u)]_5,
\]
where $b_\iota$ is multiplication by $\operatorname{res}_\iota(\beta)$; no action of
$\delta$ on $H$ is required.
\end{proposition}

\begin{proof}
Lemma~\ref{lem:cyclotomic-k4} and the coefficient exact sequence give
\[
 0\longrightarrow K_5(S)/5\longrightarrow K_5(S;5)
 \longrightarrow K_4(S)[5]=0.
\]
Consequently $K_5(S)\to K_5(S;5)$ is surjective, and
$b^2[u]_5$ has an integral lift $c$.  Applying $\delta$ after reduction,
then using naturality and bilinearity of the parenthesized product, gives
\[
 \begin{aligned}
 (\delta c)\bmod5
 &=\delta\bigl(\beta\cdot(\beta\cdot[u]_5)\bigr)\\
 &=(\delta\beta)\cdot\bigl((\delta\beta)\cdot[\delta u]_5\bigr)\\
 &=\chi(\delta)^2\beta\cdot(\beta\cdot[\delta u]_5).
 \end{aligned}
\]
This proves \eqref{eq:conjugate-c}; functoriality gives its restriction
along $\iota$.  The construction made no equivariant choice of the
integral lift $c$.
\end{proof}
